\section{History}
\begin{frame}{History}
  \centering
  
    \includegraphics[width=0.8\textwidth]{history.jpg}
\end{frame}

\begin{frame}{History}
  \begin{itemize}
  \item Regular Type Expressions (RTEs) were introduced at 2016 European Lisp Conference, implemented in Common Lisp.
    \url{https://www.lrde.epita.fr/wiki/Publications/newton.16.els}
  \item Many theoretical, implementation, and performance details explained in PhD Thesis in 2018.
    \url{https://www.lrde.epita.fr/wiki/Publications/newton.18.phd}
  \item In mean time, extended RTEs to Scala, Clojure, and Python
  \item Available:
    \begin{itemize}
    \item \url{https://github.com/jimka2001/scala-rte}
    \item \url{https://github.com/jimka2001/clojure-rte}
    \item \url{https://github.com/jimka2001/python-rte}
    \item \url{https://github.com/jimka2001/cl-rte}
    \end{itemize}
  \end{itemize}
\end{frame}


\begin{frame}{Requirements}
  \begin{itemize}
  \item Heterogeneous sequences are commonplace in dynamically typed languages
  \item Supported also in Scala via \code{Seq[Any]}
  \item Run-time type reflection 
    \begin{itemize}
    \item type membership query:\code{classOf[Int].isInstance("hello")}
    \item subtype query: \code{classOf[Int].isAssignableFrom(classOf[String])}
    \end{itemize}
  \item Implementation of a Simple Embeddable Type System
    \url{https://www.lrde.epita.fr/wiki/Publications/newton.21.els}
  \end{itemize}
\end{frame}
